\documentclass[11pt,a4paper]{article}

\usepackage{fontspec}
\usepackage{polyglossia}
\setmainlanguage{russian}
\setotherlanguage{english}
\setmainfont{PT Serif}
\setsansfont{Arial}
\setmonofont{Menlo}
\newfontfamily\cyrillicfont{PT Serif}
\newfontfamily\cyrillicfontsf{Arial}
\newfontfamily\cyrillicfonttt{Menlo}

\usepackage[a4paper,margin=19mm,headheight=14pt]{geometry}
\usepackage{microtype}
\usepackage{amsmath,amssymb}
\usepackage{booktabs}
\usepackage{tabularx}
\usepackage{longtable}
\usepackage{array}
\usepackage{enumitem}
\usepackage{xcolor}
\usepackage{colortbl}
\usepackage{fancyhdr}
\usepackage{titlesec}
\usepackage{hyperref}
\usepackage{xurl}

\definecolor{navy}{HTML}{183153}
\definecolor{blue}{HTML}{2B6CB0}
\definecolor{cyan}{HTML}{DDF2F7}
\definecolor{lightblue}{HTML}{EDF5FC}
\definecolor{lightgray}{HTML}{F3F4F6}
\definecolor{midgray}{HTML}{697386}
\definecolor{green}{HTML}{287A55}
\definecolor{red}{HTML}{B23A48}
\definecolor{gold}{HTML}{A66B00}

\hypersetup{
  colorlinks=true,
  linkcolor=navy,
  urlcolor=blue,
  citecolor=green,
  pdfauthor={Codex / GMM research project},
  pdftitle={Экспериментальный отчёт: dimension-free moment GMM},
  pdfsubject={Сводный анализ пяти экспериментов}
}

\pagestyle{fancy}
\fancyhf{}
\lhead{\sffamily\small Dimension-free moment GMM}
\rhead{\sffamily\small Экспериментальный отчёт}
\cfoot{\sffamily\small\thepage}
\renewcommand{\headrulewidth}{0.3pt}

\titleformat{\section}
  {\Large\sffamily\bfseries\color{navy}}
  {\thesection}{0.6em}{}
\titleformat{\subsection}
  {\large\sffamily\bfseries\color{blue}}
  {\thesubsection}{0.6em}{}
\titleformat{\subsubsection}
  {\normalsize\sffamily\bfseries\color{navy}}
  {\thesubsubsection}{0.6em}{}
\titlespacing*{\section}{0pt}{1.5em}{0.65em}
\titlespacing*{\subsection}{0pt}{1.15em}{0.45em}

\setlist[itemize]{leftmargin=1.3em,itemsep=0.25em,topsep=0.3em}
\setlist[enumerate]{leftmargin=1.5em,itemsep=0.3em,topsep=0.3em}
\renewcommand{\arraystretch}{1.16}
\setlength{\tabcolsep}{4.5pt}
\setlength{\parindent}{0pt}
\setlength{\parskip}{0.48em}
\emergencystretch=2em

\newcolumntype{Y}{>{\raggedright\arraybackslash}X}
\newcolumntype{R}{>{\raggedleft\arraybackslash}X}
\newcommand{\code}[1]{{\footnotesize\nolinkurl{#1}}}
\newcommand{\best}[1]{\textbf{\color{green}#1}}
\newcommand{\warn}[1]{\textbf{\color{red}#1}}
\newcommand{\note}[1]{%
  \begin{center}
  \colorbox{lightblue}{\parbox{0.94\linewidth}{#1}}
  \end{center}
}
\newcommand{\decision}[1]{%
  \begin{center}
  \colorbox{cyan}{\parbox{0.94\linewidth}{%
    \textbf{\sffamily\color{navy}Решение по результатам.} #1}}
  \end{center}
}
\newcommand{\configlink}[1]{%
  \code{#1.yaml}}
\newcommand{\resultslink}[1]{%
  \code{#1/}}
\newcommand{\runrefs}[1]{%
  \small\color{midgray}%
  Конфиг: \configlink{#1}\enspace|\enspace
  каталог результатов: \resultslink{#1}.%
  \normalsize\color{black}}

\begin{document}

\begin{titlepage}
  \pagecolor{navy}
  \color{white}
  \vspace*{24mm}
  {\sffamily\bfseries\fontsize{30}{35}\selectfont
    Dimension-free moment GMM\par}
  \vspace{7mm}
  {\sffamily\fontsize{19}{24}\selectfont
    Экспериментальный отчёт\par}
  \vspace{8mm}
  {\Large
    От глобального обзора модели до проверки устойчивости\\
    к contamination и направленной асимметрии\par}

  \vfill
  \colorbox{white}{%
    \parbox{0.86\textwidth}{%
      \color{navy}\vspace{3mm}
      \textbf{\sffamily Охват отчёта}\par
      Пять последовательных экспериментов, 74 сценария данных и пять
      взаимосвязанных семейств абляций. В отчёте сохранены только
      интерпретируемые срезы таблиц, но описаны все участвовавшие модели.
      \vspace{3mm}}}
  \vfill
  {\sffamily 30 июля 2026 г.\par}
  {\small Репозиторий: \code{/Users/boriszhukov/Projects/gmm_clean}\par}
\end{titlepage}
\nopagecolor

\tableofcontents
\clearpage

\section{Краткое резюме}

Цель серии экспериментов --- понять, какие части обобщённой моментной модели
действительно улучшают оценивание гауссовских смесей, и проверить гипотезу,
что критерий на локализованных тестовых функциях устойчивее likelihood-based
методов при малой выборке и misspecification.

\begin{center}
\small
\begin{tabularx}{\linewidth}{@{}p{0.17\linewidth}Y Y@{}}
\toprule
\rowcolor{lightgray}
\textbf{Этап} & \textbf{Главный вопрос} & \textbf{Краткий ответ}\\
\midrule
Global sweep &
Какие режимы новой модели вообще жизнеспособны? &
Стабилизация вокруг KMeans и большое число тестовых центров критичны.
EM остаётся общим лидером.\\
Clean robustness &
Что дают mean penalty, near-mean centers и вторые моменты? &
Основной вклад даёт mean penalty; near-mean centers полезны как дополнение.
Усиление вторых моментов улучшает нестабилизированную модель, но не закрывает
разрыв с EM.\\
Contamination &
Устойчив ли метод к тяжёлым хвостам и выбросам? &
EM лучше по плотности и параметрам в среднем. Moment-метод показывает
интересное преимущество по ARI на некоторых тяжёлых Student-$t$ сценариях.\\
Density centers &
Помогает ли отбрасывать разреженные центры тестовых функций? &
Иногда --- особенно для Student-$t$ и нестабильной $n_2=3$ модели.
Универсального улучшения нет; агрессивная фильтрация не оправдана.\\
Directional skewness &
Побеждает ли moment-метод EM при асимметричных компонентах? &
Нет в исследованном режиме. EM выигрывает True KL и восстановление
компонент; разрыв растёт при малой выборке и неравных весах.\\
\bottomrule
\end{tabularx}
\end{center}

\textbf{Главный содержательный вывод.}
Собранная модель уже является рабочим оценивателем: в сложных balanced
сценариях она существенно улучшает KMeans и по ARI приближается к EM.
Однако имеющиеся данные не подтверждают общее превосходство над EM по
устойчивости. Наиболее перспективная ниша, которую стоит проверять дальше,
--- асимметричные тяжёлые хвосты и gross contamination, а не просто
light-tailed skewness.

\textbf{Практический shortlist моделей.}

\begin{itemize}
  \item \code{stab_balanced_both} --- основной вариант, когда важны
        density fit и общая стабильность;
  \item \code{stab_n3_w025_near} --- наиболее разумный вариант для
        структурного восстановления компонент;
  \item \code{stab_n3_w025_both} --- дополнительный вариант для Sliced
        Wasserstein и проверки трёх вторых направлений;
  \item \code{em_kmeans_reg1e6} и \code{em_kmeans_reg1e2} ---
        обязательные likelihood baselines;
  \item \code{kmeans_spherical} --- контроль качества общей инициализации.
\end{itemize}

\section{Модель, метрики и правила чтения результатов}

\subsection{Что оптимизирует moment-модель}

Модель сопоставляет наблюдаемые ответы локализованных гауссовских тестовых
функций с их аналитическим ожиданием под сферической Gaussian mixture.
В конфигурации можно независимо выбирать количество нулевых, первых и вторых
моментов, направления второго порядка, веса семейств моментов, стратегию
работы с bandwidth $s$, геометрическую оптимизацию и ограничения на
амплитуды.

В окончательной базовой конфигурации используются:

\begin{itemize}
  \item KMeans-инициализация с десятью рестартами;
  \item автоматический выбор bandwidth;
  \item simplex-ограничение на веса смеси;
  \item совместная оптимизация всех $s$;
  \item покомпонентная геометрическая оптимизация;
  \item нормировка loss каждого порядка по числу его тестовых функций.
\end{itemize}

Два механизма стабилизации имеют разный смысл:

\begin{itemize}
  \item \textbf{mean penalty} штрафует удаление средних от KMeans-инициализации;
  \item \textbf{near-mean centers} добавляет тестовые центры около текущих
        средних и локально усиливает идентификацию компонент.
\end{itemize}

\subsection{Метрики}

Для первых четырёх экспериментов \textbf{Forward KL} вычисляется относительно
истинной гауссовской смеси или её moment-matched Gaussian core. В
contamination-экспериментах это \emph{не} KL между фактическим
негауссовским распределением наблюдений и найденной моделью.

В directional-skewness эксперименте добавлены:

\begin{itemize}
  \item \textbf{True Forward KL} --- Monte Carlo оценка
    $\mathbb E_P[\log p(X)-\log q(X)]$ для точной skewed-плотности;
  \item \textbf{Held-out NLL} на общей независимой выборке;
  \item \textbf{Sliced W2} по случайным и истинным skewness-направлениям;
  \item \textbf{Energy Distance} между обучающей $X$ и выборкой из модели;
  \item \textbf{Conditional NLL} после сопоставления истинных и найденных
    компонент.
\end{itemize}

Остальные метрики:

\begin{itemize}
  \item \textbf{Weights TV} --- total variation между векторами весов;
  \item \textbf{Means Error} --- нормированная по масштабу компонент ошибка
    средних;
  \item \textbf{Covariance Error} --- аффинно-инвариантная ошибка ковариаций;
  \item \textbf{ARI} --- permutation-invariant качество кластеризации.
\end{itemize}

\subsection{Правила агрегации и ограничения}

\begin{itemize}
  \item Во всех пяти финальных наборах каждая строка имеет \code{ok=1}.
        Строки с падениями не входят в таблицы.
  \item В отчёте не анализируется runtime: на этой стадии исследуется качество,
        а не вычислительный бюджет.
  \item Средние по разнородным сценариям используются как описательная сводка;
        основной смысл имеют сравнения внутри сценария и средние ранги.
  \item Сохранены mean/median/std/IQR, но не индивидуальные paired результаты.
        Поэтому формальные доверительные интервалы для разности методов
        восстановить нельзя.
  \item Три model seeds на одном dataset seed не являются тремя независимыми
        наборами данных. Стандартные отклонения в CSV нельзя напрямую
        интерпретировать как стандартную ошибку среднего.
  \item Выбор ``лучшей moment-модели'' после просмотра результатов является
        оптимистичным; для подтверждающего эксперимента shortlist должен быть
        зафиксирован заранее.
\end{itemize}

\section{Каталог всех участвовавших моделей}

\subsection{Варианты первого global sweep}

\small
\begin{longtable}{@{}>{\raggedright\arraybackslash}p{0.42\linewidth}
                         >{\raggedright\arraybackslash}p{0.52\linewidth}@{}}
\toprule
\rowcolor{lightgray}
\textbf{Имя модели} & \textbf{Изменение относительно базовой moment-конфигурации}\\
\midrule
\endfirsthead
\toprule
\rowcolor{lightgray}
\textbf{Имя модели} & \textbf{Изменение относительно базовой moment-конфигурации}\\
\midrule
\endhead
\code{moment_first_only3} & Только три первых направления: $(n_0,n_1,n_2)=(0,3,0)$.\\
\code{moment_second_only3} & Только три вторых направления: $(0,0,3)$.\\
\code{moment_zero1_first2} & Нулевой и два первых момента: $(1,2,0)$.\\
\code{moment_first2_second1} & Два первых и один второй момент: $(0,2,1)$.\\
\code{moment_zero1_second2} & Нулевой и два вторых момента: $(1,0,2)$.\\
\code{moment_all_balanced} & База $(1,1,1)$ с равными весами семейств.\\
\code{moment_all_first_rich} & $(1,3,1)$: больше первых направлений.\\
\code{moment_all_second_rich} & $(1,1,3)$: больше вторых направлений.\\
\code{moment_all_cross_second} & $(1,1,1)$, второй момент использует два независимых направления.\\
\code{moment_all_second_weight025} & База с весом второго семейства $0.25$.\\
\code{moment_all_second_weight4} & База с весом второго семейства $4.0$.\\
\code{moment_all_non_negative} & Неотрицательные амплитуды без simplex-связи во время оптимизации.\\
\code{moment_all_sequential_s} & Последовательная оптимизация по bandwidth $s$.\\
\code{moment_all_joint_geometry} & Одновременная геометрическая оптимизация всех компонент.\\
\code{moment_all_centers50_data} & 50 центров, выбранных непосредственно из данных.\\
\code{moment_all_centers50_noisy_auto} & 50 noisy-data центров; шум выбирается автоматически.\\
\code{moment_all_centers200_noisy_auto} & То же, но 200 центров.\\
\code{moment_all_stabilized} & База + near-mean centers с ratio $1.0$ + mean penalty $10^{-3}$.\\
\code{sklearn_gmm_kmeans} & Сферическая sklearn GMM/EM, KMeans init, \code{reg_covar=1e-6}.\\
\code{sklearn_kmeans} & Сферический KMeans-контроль с эмпирическими весами и дисперсиями.\\
\bottomrule
\end{longtable}
\normalsize

\subsection{Общий каталог последующих абляций}

\begin{center}
\small
\begin{tabularx}{\linewidth}{@{}>{\raggedright\arraybackslash}p{0.42\linewidth}Y@{}}
\toprule
\rowcolor{lightgray}
\textbf{Имена} & \textbf{Смысл}\\
\midrule
\code{stab_balanced_none}, \code{..._penalty}, \code{..._near},
\code{..._both} &
Факториал для $(n_0,n_1,n_2)=(1,1,1)$:
без стабилизации; только mean penalty; только near-mean centers;
оба механизма.\\
\code{second_n1_w010} \ldots \code{second_n5_w050} &
Полная сетка $n_2\in\{1,3,5\}$ и
$w_2\in\{0.10,0.25,0.50\}$ при одном нулевом и одном первом моменте.\\
\code{second_n3_w025} и
\code{stab_n3_w025_penalty/near/both} &
Повторный стабилизационный факториал в выбранной точке $n_2=3,w_2=0.25$.\\
\code{second_n5_w050} &
Сильная граница second-moment grid, использованная в contamination-тестах.\\
\code{moment_first_only3} &
Контроль без нулевых и вторых моментов.\\
\code{em_kmeans_reg1e6}, \code{em_kmeans_reg1e2} &
Сферическая EM с KMeans init и регуляризацией ковариации
$10^{-6}$ или $10^{-2}$.\\
\code{kmeans_spherical} &
Инициализационный контроль без последующей EM/moment-оптимизации.\\
\bottomrule
\end{tabularx}
\end{center}

В clean robustness эксперименте участвовали все четыре balanced-варианта,
все девять точек second-moment grid, три стабилизированных варианта
$n_2=3,w_2=0.25$, first-only контроль и три baseline --- всего 20 моделей.

В contamination и directional-skewness экспериментах использован сокращённый
набор из десяти moment-моделей:
\code{stab_balanced_none/penalty/near/both},
\code{second_n3_w025},
\code{stab_n3_w025_penalty/near/both},
\code{second_n5_w050},
\code{moment_first_only3};
плюс две EM и KMeans --- всего 13.

В density-center эксперименте использованы пять семейств:
\code{stab_balanced_both}, \code{stab_balanced_penalty},
\code{stab_balanced_near}, \code{stab_n3_w025_both},
\code{stab_balanced_none}. Каждое семейство протестировано с
\code{test_center_fraction} $f\in\{1.0,0.9,0.5,0.3\}$; суффиксы
\code{f100/f090/f050/f030} однозначно описывают все 20 moment-моделей.
Дополнительно включены те же три baseline.

\section{Эксперимент 1: глобальный обзор модели}

\runrefs{optimization_moments_global_sweep}

\subsection{Постановка}

Эксперимент охватывает 20 сценариев: решётку
$d\in\{3,30\}$, $n\in\{60,500\}$ и easy/hard separation;
стресс по дисперсиям и весам; смеси из пяти компонент; диагональную и полную
ковариации при сферической оцениваемой модели. Для каждого сценария использовано 10 dataset
seeds и 2 model seeds. Всего сравнивались 18 moment-вариантов и две sklearn
модели.

\subsection{Сокращённая таблица результатов}

\begin{center}
\small
\begin{tabular}{@{}lrrrrr@{}}
\toprule
\rowcolor{lightgray}
\textbf{Модель} & \textbf{KL} $\downarrow$ & \textbf{TV} $\downarrow$ &
\textbf{Mean} $\downarrow$ & \textbf{Cov} $\downarrow$ & \textbf{ARI} $\uparrow$\\
\midrule
EM & \best{0.325} & \best{0.132} & \best{0.220} & \best{0.203} & \best{0.641}\\
KMeans & 0.404 & 0.163 & 0.253 & 0.229 & 0.585\\
all stabilized & 0.391 & 0.159 & 0.240 & 0.209 & 0.600\\
first only 3 & 0.407 & 0.165 & 0.243 & 0.216 & 0.589\\
all second rich & 0.517 & 0.163 & 0.261 & 0.226 & 0.592\\
all balanced & 0.695 & 0.164 & 0.275 & 0.267 & 0.585\\
\bottomrule
\end{tabular}
\end{center}

Это средние по всем 20 сценариям; они нужны для общей ориентации, а не для
формального рейтинга. По среднему внутрисценарному рангу EM также первая,
\code{moment_all_stabilized} --- лучшая moment-модель.

EM выиграла Forward KL в 13 из 20 сценариев; first-only moment ---
в четырёх. Stabilized moment уступает EM, но улучшает исходную balanced
модель по KL на 44\%, по mean error на 13\% и по ARI примерно на 0.016.

\subsection{Что выяснили отдельные абляции}

\begin{itemize}
  \item \textbf{Стабилизация важнее выбора порядка моментов.}
        \code{all_stabilized} имеет KL $0.391$ против $0.695$ у
        \code{all_balanced}.
  \item \textbf{Первые моменты особенно информативны.}
        First-only даёт KL $0.407$ и mean error $0.243$.
  \item \textbf{Дополнительные вторые направления полезнее одного:}
        second-rich улучшает balanced KL с $0.695$ до $0.517$.
  \item \textbf{Два независимых second-order направления} дают умеренное
        улучшение KL ($0.612$), но не меняют общий вывод.
  \item \textbf{Слишком большой second weight вреден:}
        $w_2=4$ даёт KL $0.730$, а $w_2=0.25$ --- $0.644$.
  \item \textbf{Joint geometry нестабилен:} mean error $1.074$ против
        $0.275$ у покомпонентной оптимизации.
  \item \textbf{Последовательная оптимизация по $s$ уступает совместной:}
        KL $0.906$ против $0.695$.
  \item \textbf{Число центров существенно:} 50 центров дают KL
        $0.73$--$0.86$, 200 noisy centers --- $0.499$.
  \item \textbf{Non-negative амплитуды не дают устойчивого преимущества}
        и теряют прямую связь с нормированной смесью во время оптимизации;
        simplex остаётся предпочтительным дефолтом.
\end{itemize}

\decision{Зафиксировать simplex, joint-$s$, component geometry и достаточно
большое число тестовых центров. Разделить стабилизацию на mean penalty и
near-mean centers и отдельно исследовать количество/вес вторых моментов.}

\section{Эксперимент 2: clean robustness и moment ablation}

\runrefs{moments_robustness_ablation}

\subsection{Постановка}

18 чистых гауссовских сценариев включают sample-size ladder
$n\in\{30,60,200,500\}$, dimension stress до $d=60$, стресс по дисперсиям и
весам, separation ladder, пять компонент, несферические ковариации и easy
calibration. Использовано 30 dataset seeds и 3 model seeds.

Эксперимент отвечает на два заранее поставленных вопроса:

\begin{enumerate}
  \item разделить эффекты mean penalty и near-mean centers;
  \item проверить сетку $n_2\in\{1,3,5\}$,
        $w_2\in\{0.10,0.25,0.50\}$.
\end{enumerate}

\subsection{Факториал стабилизации}

\begin{center}
\small
\begin{tabular}{@{}lrrrrr@{}}
\toprule
\rowcolor{lightgray}
\textbf{Модель} & \textbf{KL} $\downarrow$ & \textbf{TV} $\downarrow$ &
\textbf{Mean} $\downarrow$ & \textbf{Cov} $\downarrow$ & \textbf{ARI} $\uparrow$\\
\midrule
balanced none & 1.972 & 0.215 & 0.448 & 1.722 & 0.499\\
balanced penalty & 0.937 & 0.205 & 0.332 & \best{0.674} & 0.557\\
balanced near & 1.851 & 0.208 & 0.414 & 1.324 & 0.511\\
balanced both & \best{0.888} & \best{0.204} & \best{0.329} & 0.704 & \best{0.560}\\
\midrule
$n_2=3$ none & 1.654 & 0.222 & 0.444 & 1.419 & 0.500\\
$n_2=3$ penalty & 0.972 & \best{0.212} & 0.333 & \best{0.742} & 0.547\\
$n_2=3$ near & 1.500 & 0.221 & 0.383 & 1.038 & 0.512\\
$n_2=3$ both & \best{0.968} & 0.213 & \best{0.331} & 0.773 & \best{0.547}\\
\bottomrule
\end{tabular}
\end{center}

Mean penalty является главным стабилизирующим фактором. Near-mean centers
сами по себе улучшают исходную модель гораздо слабее, но вместе со штрафом
дают небольшое дополнительное улучшение KL, весов, средних и ARI.
У $n_2=3$ добавление near centers поверх penalty почти нейтрально по ARI.

\subsection{Сетка вторых моментов}

\begin{center}
\small
\begin{tabular}{@{}ccrr@{\qquad}ccrr@{}}
\toprule
\rowcolor{lightgray}
$n_2$ & $w_2$ & \textbf{KL} $\downarrow$ & \textbf{ARI} $\uparrow$ &
$n_2$ & $w_2$ & \textbf{KL} $\downarrow$ & \textbf{ARI} $\uparrow$\\
\midrule
1 & 0.10 & 1.830 & 0.481 & 3 & 0.10 & 1.674 & 0.490\\
1 & 0.25 & 1.886 & 0.488 & 3 & 0.25 & 1.654 & 0.500\\
1 & 0.50 & 1.926 & 0.494 & 3 & 0.50 & 1.637 & 0.508\\
\midrule
5 & 0.10 & 1.578 & 0.496 & 5 & 0.25 & 1.517 & 0.508\\
5 & 0.50 & \best{1.455} & \best{0.518} & & & &\\
\bottomrule
\end{tabular}
\end{center}

В нестабилизированной сетке больше направлений последовательно улучшает
результат. При $n_2=3$ и $5$ больший вес второго семейства также полезен.
Однако даже лучшая точка $n_2=5,w_2=0.5$ заметно хуже стабилизированной
balanced-модели. Следовательно, вторые моменты содержат полезный сигнал, но
не заменяют регуляризацию геометрии.

\subsection{Sample-size ladder}

\begin{center}
\small
\begin{tabular}{@{}r rr rr@{}}
\toprule
\rowcolor{lightgray}
& \multicolumn{2}{c}{\textbf{EM reg $10^{-2}$}} &
\multicolumn{2}{c}{\textbf{balanced both}}\\
\cmidrule(lr){2-3}\cmidrule(l){4-5}
$n$ & \textbf{KL} $\downarrow$ & \textbf{ARI} $\uparrow$ &
\textbf{KL} $\downarrow$ & \textbf{ARI} $\uparrow$\\
\midrule
30  & \best{1.194} & \best{0.477} & 1.514 & 0.471\\
60  & \best{0.495} & \best{0.644} & 0.718 & 0.570\\
200 & \best{0.091} & \best{0.780} & 0.246 & 0.694\\
500 & \best{0.035} & \best{0.788} & 0.147 & 0.717\\
\bottomrule
\end{tabular}
\end{center}

Гипотеза о более мягкой деградации moment-метода при малом $n$ не
подтвердилась. При $n=30$ методы близки по ARI, но с ростом данных EM
использует дополнительную информацию значительно эффективнее. Dimension
stress усиливает этот вывод: при $d=n=60$ EM имеет KL $4.14$ и ARI $0.468$,
balanced-both --- KL $5.95$ и ARI $0.380$.

При этом moment-оптимизация существенно лучше простого KMeans в большинстве
сложных сценариев. Это подтверждает, что критерий обучает модель, а не только
наследует сильную инициализацию.

\decision{Основным вариантом оставить \code{stab_balanced_both}; сохранить
\code{stab_n3_w025_both/near} как структурные альтернативы. Гипотезу
превосходства над EM перенести из чистого small-$n$ режима в явное
misspecification и contamination.}

\section{Эксперимент 3: contamination robustness}

\runrefs{moments_contamination_robustness}

\subsection{Постановка}

12 сценариев образованы четырьмя типами отклонения от Gaussian mixture:

\begin{itemize}
  \item Student-$t$ компоненты с $\nu\in\{3,5,10\}$;
  \item component-wise variance contamination;
  \item diffuse uniform background;
  \item удалённые radial outliers.
\end{itemize}

Для каждого сценария использовано 20 dataset seeds и 3 model seeds.
Parameter metrics и старая Forward KL направлены на чистую Gaussian core.
Для явных outliers ARI считается только по clean observations; у Student-$t$
метки определены для всех наблюдений.

\subsection{Средние по типам contamination}

\begin{center}
\small
\begin{tabular}{@{}l rr rr@{}}
\toprule
\rowcolor{lightgray}
& \multicolumn{2}{c}{\textbf{EM reg $10^{-2}$}} &
\multicolumn{2}{c}{\textbf{balanced both}}\\
\cmidrule(lr){2-3}\cmidrule(l){4-5}
\textbf{Тип} & \textbf{KL} $\downarrow$ & \textbf{ARI} $\uparrow$ &
\textbf{KL} $\downarrow$ & \textbf{ARI} $\uparrow$\\
\midrule
Radial       & \best{0.491} & \best{0.502} & 0.651 & 0.423\\
Student-$t$  & \best{0.733} & 0.442 & 1.108 & \best{0.500}\\
Uniform      & \best{0.722} & \best{0.465} & 1.082 & 0.403\\
Variance     & \best{0.630} & \best{0.397} & 0.766 & 0.348\\
\bottomrule
\end{tabular}
\end{center}

EM сохраняет преимущество по Gaussian-core density и параметрам почти во
всех группах. Moment-модель не демонстрирует общей B-robustness: удалённые
точки всё ещё влияют на выбранные средние и дисперсии через empirical moment
responses и KMeans init.

Однако Student-$t$ даёт важное исключение: по ARI moment-метод в среднем
лучше EM. Особенно заметен сценарий $\nu=3$:

\begin{center}
\small
\begin{tabular}{@{}lrr@{}}
\toprule
\rowcolor{lightgray}
\textbf{Сценарий} & \textbf{Лучший moment ARI} & \textbf{Лучший EM ARI}\\
\midrule
Student-$t$, $\nu=10$, $d=10$ & 0.670 & 0.648\\
Student-$t$, $\nu=3$, $d=10$ & \best{0.533} & 0.342\\
Student-$t$, $\nu=5$, $d=30$ & 0.309 & \best{0.336}\\
Radial, 10\%, radius 10--16 & \best{0.129} & 0.079\\
Uniform, 10\%, padding 8 & \best{0.115} & 0.063\\
Variance, 10\%, inflation 100 & \best{0.063} & 0.046\\
\bottomrule
\end{tabular}
\end{center}

Последние три преимущества достигаются при очень низком абсолютном ARI:
это скорее более медленное разрушение, чем практически хорошая
кластеризация. А вот разница на Student-$t$ с $\nu=3$ содержательна и
поддерживает гипотезу о потенциальном преимуществе bounded/localized
моментов при тяжёлых хвостах.

Из moment-моделей наиболее устойчив \code{stab_balanced_both}.
Mean penalty ограничивает уход средних под влиянием выбросов; near centers
без penalty иногда ухудшают mean error, поскольку текущая геометрия уже могла
сместиться к contamination. Сырые $n_2=3/5$ варианты не выигрывают.

\decision{Не заявлять общую contamination-robustness. Зафиксировать
Student-$t$ с малым $\nu$ как главный положительный сигнал. Проверить,
можно ли уменьшить влияние outliers через выбор только плотных data-centered
тестовых функций.}

\section{Эксперимент 4: density-based выбор тестовых центров}

\runrefs{moments_density_center_ablation}

\subsection{Постановка}

Для каждой точки данных вычисляется ``видимость''
$\sum_i K_s(X_i,c)$; центрами тестовых функций становятся только точки из
верхней доли $f\in\{1.0,0.9,0.5,0.3\}$. Использованы те же 12 contaminated
сценариев, 8 dataset seeds и 3 model seeds. Пять moment-семейств образуют
20 вариантов; EM и KMeans сохранены для прямого сравнения.

\subsection{Средний эффект доли центров}

\begin{center}
\small
\begin{tabular}{@{}crrrrr@{}}
\toprule
\rowcolor{lightgray}
$f$ & \textbf{KL} $\downarrow$ & \textbf{TV} $\downarrow$ &
\textbf{Mean} $\downarrow$ & \textbf{Cov} $\downarrow$ & \textbf{ARI} $\uparrow$\\
\midrule
1.0 & 1.058 & 0.227 & 2.105 & 1.828 & 0.394\\
0.9 & \best{0.997} & 0.225 & \best{2.040} & 1.841 & 0.389\\
0.5 & 0.998 & 0.221 & 2.043 & 1.818 & 0.393\\
0.3 & 1.010 & \best{0.217} & 2.052 & \best{1.813} & \best{0.395}\\
\bottomrule
\end{tabular}
\end{center}

Пул по всем семействам показывает умеренное улучшение KL и весов, но не
монотонный эффект и почти отсутствие выигрыша по ARI. Пул несколько
искажён тем, что фильтрация сильно спасает отдельные нестабильные семейства.

Для наиболее сильного \code{stab_balanced_both}:

\begin{center}
\small
\begin{tabular}{@{}crrrrr@{}}
\toprule
\rowcolor{lightgray}
$f$ & \textbf{KL} $\downarrow$ & \textbf{TV} $\downarrow$ &
\textbf{Mean} $\downarrow$ & \textbf{Cov} $\downarrow$ & \textbf{ARI} $\uparrow$\\
\midrule
1.0 & \best{0.867} & 0.223 & \best{1.732} & 1.862 & 0.414\\
0.9 & 0.908 & 0.218 & 1.734 & 1.882 & 0.404\\
0.5 & 0.889 & 0.214 & 1.749 & \best{1.851} & 0.413\\
0.3 & 0.882 & \best{0.212} & 1.740 & 1.863 & \best{0.414}\\
\bottomrule
\end{tabular}
\end{center}

Здесь unfiltered $f=1$ сохраняет лучшую KL и mean error. $f=0.3$ улучшает
веса без потери среднего ARI, но не даёт общего доминирования.

\subsection{Зависимость от типа contamination}

По pooled Forward KL:

\begin{center}
\small
\begin{tabular}{@{}lrrrr@{}}
\toprule
\rowcolor{lightgray}
\textbf{Тип} & $f=1.0$ & $f=0.9$ & $f=0.5$ & $f=0.3$\\
\midrule
Radial      & \best{0.663} & 0.683 & 0.701 & 0.704\\
Student-$t$ & 1.586 & \best{1.272} & 1.277 & 1.304\\
Uniform     & 1.251 & 1.257 & \best{1.226} & 1.245\\
Variance    & \best{0.731} & 0.776 & 0.790 & 0.788\\
\bottomrule
\end{tabular}
\end{center}

Фильтр особенно полезен для Student-$t$: он снижает влияние периферийных
наблюдений как тестовых центров. Для radial и variance contamination
удалённые центры, напротив, несут полезную информацию о несогласованности
модели с данными; их исключение ухудшает KL к clean core в выбранной
оптимизационной динамике. Для uniform background эффект слаб и неоднозначен.

EM остаётся лидером общего рейтинга; density filtering не создаёт
преимущества moment-метода над likelihood baselines.

\decision{Не менять дефолт $f=1.0$. Рассматривать $f=0.9$ или $0.5$ как
специализированную robust-настройку для тяжёлых хвостов; $f=0.3$ слишком
агрессивен для общего применения. В будущем выбирать $f$ адаптивно по
распределению density scores, а не фиксированной долей.}

\section{Эксперимент 5: направленная асимметрия компонент}

\runrefs{moments_directional_skewness}

\subsection{Постановка}

Directional-Gamma генератор сохраняет точные истинные средние и ковариации.
В whitened координатах одна компонента residual заменяется на
\[
S=\frac{G-a}{\sqrt a},\qquad G\sim\Gamma(a,1),
\]
а остальные направления остаются гауссовскими. При $a=2$ стандартизованная
skewness равна $2/\sqrt{2}$. Направление длинного хвоста выбирается как:

\begin{itemize}
  \item \code{random} --- случайное для каждой компоненты;
  \item \code{toward} --- к барицентру смеси;
  \item \code{away} --- от барицентра.
\end{itemize}

Два data regimes (balanced $n=500$ и scarce/moderately imbalanced $n=200$)
пересечены с separation 3.5/6.0 и тремя направлениями --- всего 12
сценариев. Использовано 20 dataset seeds, 3 model seeds и независимая
evaluation sample размера 10\,000.

\subsection{Общий результат}

\begin{center}
\small
\begin{tabular}{@{}lrrrrrr@{}}
\toprule
\rowcolor{lightgray}
\textbf{Модель} & \textbf{True KL} $\downarrow$ & \textbf{SW2} $\downarrow$ &
\textbf{TV} $\downarrow$ & \textbf{Mean} $\downarrow$ &
\textbf{Cov} $\downarrow$ & \textbf{ARI} $\uparrow$\\
\midrule
EM reg $10^{-2}$ & \best{0.242} & \best{0.147} & \best{0.035} &
0.109 & \best{0.049} & \best{0.962}\\
EM reg $10^{-6}$ & 0.244 & 0.147 & 0.036 &
\best{0.108} & 0.052 & 0.961\\
Лучший moment по каждой метрике & 0.283 & 0.153 & 0.050 &
0.132 & 0.084 & 0.942\\
KMeans & 0.323 & 0.153 & 0.063 &
0.140 & 0.144 & 0.928\\
\bottomrule
\end{tabular}
\end{center}

EM reg $10^{-2}$ имеет лучшую mean True KL во всех 12 сценариях.
EM также выигрывает Conditional NLL в 11 из 12 и почти все parameter
metrics. Energy Distance практически не различает методы
($0.1582$--$0.1587$ у лидеров) и остаётся вторичной диагностикой.

\subsection{Взаимодействие sample regime и separation}

Для интерпретации дополнительно вычислена oracle KL Gaussian mixture с
точными истинными весами, средними и ковариациями. Это не новое обучение, а
оценка неизбежной ошибки Gaussian approximation.

\begin{center}
\small
\begin{tabular}{@{}lrrrrr@{}}
\toprule
\rowcolor{lightgray}
\textbf{Режим} & \textbf{Oracle KL} &
\textbf{EM KL} & \textbf{$n_2=3$ near KL} &
\textbf{EM ARI} & \textbf{Moment ARI}\\
\midrule
balanced easy & 0.184 & \best{0.216} & 0.221 & 0.993 & 0.993\\
balanced hard & 0.171 & \best{0.202} & 0.222 & 0.947 & 0.945\\
scarce easy   & 0.184 & \best{0.277} & 0.326 & 0.987 & 0.962\\
scarce hard   & 0.173 & \best{0.275} & 0.376 & 0.921 & 0.863\\
\bottomrule
\end{tabular}
\end{center}

При balanced-easy методы почти эквивалентны. Разрыв растёт при overlap и
особенно при дефиците данных/неравных весах. В scarce-hard estimation excess
над oracle составляет около $0.10$ для EM и $0.20$ для выбранной
moment-модели. Предполагаемое преимущество moment-метода в failure tails
не наблюдается: его variance и разрыв mean--median часто выше.

\subsection{Эффект направления}

В balanced-hard постановке направление работает ожидаемо:

\begin{center}
\small
\begin{tabular}{@{}lrr@{}}
\toprule
\rowcolor{lightgray}
\textbf{Направление} & \textbf{EM ARI} & \textbf{$n_2=3$ near ARI}\\
\midrule
away   & 0.980 & 0.980\\
random & 0.947 & 0.943\\
toward & 0.914 & 0.912\\
\bottomrule
\end{tabular}
\end{center}

При \code{toward} компоненты сильнее перекрываются. Одновременно True KL
становится ниже: marginal density при inward tails легче приблизить смесью,
даже если latent labels хуже идентифицируются. Это наглядно показывает,
почему density fit и recovery компонент нельзя сводить к одной метрике.

В scarce-imbalanced режиме \code{away} тоже становится сложным: внешний
хвост большой компоненты способен сместить оценивание и скрыть малую
компоненту. Направление взаимодействует с весами, а не задаёт универсальный
порядок сложности.

\subsection{Статистические вырождения}

Conditional NLL выявила редкие случаи, где \code{ok=1}, median score около
14, но mean достигает $10^6$--$10^7$. Это соответствует содержательно
неудачному решению: малая истинная компонента пропущена или сопоставлена с
почти вырожденной найденной компонентой. Проблема встречается в основном у
нестабилизированных moment-вариантов. Mean penalty устраняет наиболее
экстремальные случаи.

\decision{Гипотеза о преимуществе над EM при light-tailed directional
skewness не подтверждена. Сохранить \code{stab_balanced_both} и
\code{stab_n3_w025_near}; следующий тест строить как dose-response по силе
skewness и как skew-$t$, где асимметрия сочетается с тяжёлым хвостом.}

\section{Сквозные выводы о работе метода}

\subsection{Что подтверждено}

\begin{enumerate}
  \item \textbf{Модель действительно обучается и улучшает KMeans.}
        В balanced-hard skewness сценариях KMeans имеет KL $0.315$, moment
        около $0.222$, EM $0.202$; ARI moment почти совпадает с EM.
  \item \textbf{Simplex --- правильный дефолт.}
        Он сохраняет интерпретацию реальной смеси и не уступает
        non-negative варианту систематически.
  \item \textbf{Component geometry стабильнее joint geometry.}
        Совместная геометрия слишком чувствительна и даёт крупные ошибки
        средних.
  \item \textbf{Mean penalty --- основной стабилизатор.}
        Он снижает KL, parameter error и частоту статистических вырождений.
  \item \textbf{Near-mean centers полезны как дополнение.}
        Они улучшают локальную идентификацию, веса и иногда ARI, но без
        anchoring могут следовать уже смещённой геометрии.
  \item \textbf{Вторые моменты полезны, но требуют баланса.}
        Больше направлений улучшает сырую модель; выигрыш меньше эффекта
        стабилизации. Умеренный вес $0.25$ остаётся разумным общим выбором.
  \item \textbf{Много тестовых центров лучше малого фиксированного бюджета.}
        Это было явно видно уже в global sweep.
  \item \textbf{Density filtering не является бесплатной robustness.}
        Реальные хвостовые точки могут нести полезный сигнал; фильтр следует
        применять только под конкретную contamination model.
\end{enumerate}

\subsection{Что не подтверждено}

\begin{itemize}
  \item Нет общего преимущества над EM при малой чистой выборке.
  \item Нет общего преимущества по Gaussian-core KL при contamination.
  \item Нет преимущества при направленной Gamma-skewness.
  \item Увеличение числа вторых моментов само по себе не решает проблемы
        локальной оптимизации и component stealing.
  \item Energy Distance на training $X$ недостаточно чувствительна для
        ранжирования близких моделей.
\end{itemize}

\subsection{Где остаётся положительная гипотеза}

Наиболее сильный сигнал --- ARI на Student-$t$ с $\nu=3$:
moment $0.533$ против EM $0.342$. Это согласуется с идеей, что
локализованные kernel moments могут снижать влияние тяжёлых хвостов.
Directional Gamma имеет экспоненциальный, а не полиномиальный хвост и не
проверяет этот механизм в полной мере.

\section{Рекомендуемая следующая программа экспериментов}

\begin{enumerate}
  \item \textbf{Зафиксировать shortlist} до запуска:
    \code{stab_balanced_both}, \code{stab_n3_w025_near},
    \code{stab_n3_w025_both}, две EM и KMeans.
  \item \textbf{Skewness dose-response:}
    $\texttt{gamma\_shape}\in\{0.5,1,2,8\}$ плюс Gaussian calibration.
  \item \textbf{Skew-$t$:} независимо менять degrees of freedom и силу
    skewness. Это прямой тест наиболее перспективной гипотезы.
  \item \textbf{Adaptive density filtering:} выбирать центры по квантилям
    density score только при диагностированном тяжёлом хвосте; сравнить
    $f=1$, $0.9$, $0.5$ без $0.3$.
  \item \textbf{Failure diagnostics:} сохранять minimum weight, minimum
    variance, maximum matching distance и долю Conditional NLL выше порога.
  \item \textbf{Сохранять raw paired metrics,} включая dataset seed и model
    seed. Это позволит строить paired bootstrap intervals и разделить
    вариативность данных от вариативности оптимизации.
  \item \textbf{Добавить Gaussian oracle строкой в результаты,} чтобы
    отделять approximation error от estimation error без отдельного
    постанализа.
  \item \textbf{Развести две конечные цели:}
    density approximation оценивать True KL/NLL/SW2, а recovery компонент ---
    ARI, parameter errors и Conditional NLL. Не объединять их в один score.
\end{enumerate}

\section{Карта артефактов}

\begin{center}
\small
\begin{tabularx}{\linewidth}{@{}p{0.25\linewidth}Y Y@{}}
\toprule
\rowcolor{lightgray}
\textbf{Эксперимент} & \textbf{Конфигурация} & \textbf{Результаты}\\
\midrule
Global sweep &
\configlink{optimization_moments_global_sweep} &
\resultslink{optimization_moments_global_sweep}\\
Clean robustness &
\configlink{moments_robustness_ablation} &
\resultslink{moments_robustness_ablation}\\
Contamination &
\configlink{moments_contamination_robustness} &
\resultslink{moments_contamination_robustness}\\
Density centers &
\configlink{moments_density_center_ablation} &
\resultslink{moments_density_center_ablation}\\
Directional skewness &
\configlink{moments_directional_skewness} &
\resultslink{moments_directional_skewness}\\
\bottomrule
\end{tabularx}
\end{center}

\vfill
\begin{center}
\small\color{midgray}
Отчёт построен по финальным CSV в каталоге \code{results/}.\\
Ни один эксперимент при подготовке отчёта не перезапускался.
\end{center}

\end{document}
